\label{ch:conclusions}

In this dissertation, computational method of neutron noise is introduced as a framework to detect anomalies in the HTGR core. This dissertation has two primary goals. First, we would like to introduce and apply the computational method of neutron noise analysis to a prismatic-type HTGR core, utilizing existing noise models to determine if neutron noise analysis can be generally applied to detect anomalies in advanced reactors, primarily prismatic-type HTGRs. The second goal is to two new methods of neutron noise unfolding. These methods are developed such that it can unfold multiple noise sources simultaneously. This is achieved by utilizing the Green's function, which has been employed in existing neutron noise unfolding methods.

The following sections provide a summary of the main findings, discuss how the research objectives have been addressed, and highlight possible directions for future work.

\section{Summary of Findings}

A literature review of neutron noise analysis is provided in Chapter \ref{ch:literature_review}. In this chapter, the theory and application of neutron noise are highlighted. The goal of this chapter is to provide a summary of past neutron noise analysis, which motivates the development of computational methods for neutron noise analysis. This chapter presents the derivation of the neutron noise equation from the multigroup diffusion equation in the frequency domain, as well as the derivation of other parameters crucial for neutron noise analysis. This includes the derivation of the point-kinetic and spatial components of the neutron noise, as well as the zero-power reactor transfer function. This chapter also provides examples of solvers developed to solve the neutron noise equation, utilizing both Monte Carlo and deterministic methods.

In Chapter \ref{ch:simulations}, a numerical solver to solve the forward, adjoint, and neutron noise equations is developed. This solver is capable of generating the Green's function matrix that will be used for the neutron noise unfolding methods. This solver is developed to fill the gap in the literature, particularly in terms of generating Green's functions in triangular geometries. The solver relies on a multigroup diffusion equation for energy and angular treatment, as well as a box-scheme finite difference method for spatial discretization. Sparse array matrices are used in the solver, which is solved using the PETSc Library. It is available in \href{https://github.com/harunardi/FANGS-UNFOLD}{GitHub} for future references. The solver has been verified with various benchmarks for both rectangular and triangular geometries. The results show that the numerical solver agrees well with numerical benchmarks and analytical solutions. This is true for the forward, adjoint, and neutron noise solutions.

In Chapter \ref{ch:noise_HTTR}, the chapter is divided into two main parts. The first part involves determining the characteristics of neutron noise in the HTTR. This is done by performing a neutron noise simulation to obtain the point-kinetic and spatial components of the neutron noise in HTTR, as well as by calculating the zero-power reactor transfer function (ZPRTF). Compared to LWRs, the neutron noise in HTTR generally follows the expected behavior of ZPRTF. A shorter plateau is observed in the ZPRTF of the HTTR case compared to other types of reactors. This means that the HTTR is more sensitive to perturbations at certain frequencies compared to other reactors. Using these results, we demonstrate that neutron noise analysis can be effectively utilized in the HTTR for core diagnostics. The second part of Chapter \ref{ch:noise_HTTR} is the application of neutron noise models in the HTTR core. The results show that for the AVS-type source, the neutron noise is more localized in the thermal group and more dispersed in the fast group. It is also observed that the point-kinetic component of the neutron noise is more dominant than the spatial component of neutron noise. For the fuel assembly vibration source, the results indicate that neutron noise has a significant impact on the overall vibration. This is more apparent when the neighboring assembly of the vibrating has a different characteristic (e.g., the neighboring assembly of the fuel assembly is a reflector assembly). 

In Chapter \ref{ch:unfolding_methods}, we introduced novel neutron noise unfolding methods, namely the brute force method and the greedy method. These methods are applied to unfold multiple AVS-type sources in 2D and 3D HTTR cores. As comparison, we also performed neutron noise unfolding using the existing method, namely the inversion method, zoning method, and scanning method. The results show that the inversion and zoning methods are not reliable for unfolding neutron noise in HTTR. The scanning method can unfold one AVS-type source, but fails to unfold two or three sources. The brute force and greedy methods can unfold the neutron noise in the 2D and 3D HTTR core. This demonstrates the capability of the novel method to unfold multiple neutron noise sources simultaneously.

\section{Future Works}

We have demonstrated how the neutron noise analysis can be applied to detect anomalies in a prismatic-type HTGR core. In doing so, we have shown that the neutron noise analysis can be performed in the HTGR core. We also introduced two novel methods for unfolding neutron noise in this study. We viewed this study as a starting point for exploring various topics in neutron noise analysis, particularly in the context of advanced reactor applications. For future work, we divide it into three categories:
\begin{itemize}
    \item For the numerical solver, we would argue that the use of the multigroup neutron diffusion equation is an industry standard. However, for spatial discretization, the nodal method is preferable in an industry setting. Therefore, the solver can be refactored to utilize the nodal method (e.g., polynomial expansion method) for the spatial discretization.
    \item For the application of neutron noise in the HTGR core, several engineering issues can be modeled and simulated using neutron noise analysis. One of those issues is the effect of flow-induced vibrations on the neutron economy and reactivity fluctuations. We argue that understanding flow-induced vibration through neutron noise analysis would be beneficial for detecting anomalies caused by vibration and assessing its impact on reactor operation.
    \item For the neutron noise unfolding methods, we can expand the utilization of the novel methods by optimizing the algorithm to obtain the correct combination of Green's functions. On the other hand, we can explore other algorithms to optimize the search algorithm. Pathfinding algorithms, such as Dijkstra's and A*, can be a good starting point in this case.
\end{itemize}